\chapter{Choroidal Melanoma}
\index{Choroidal Melanoma}

\section*{Definition and Overview}
Choroidal melanoma is the most common primary malignant tumour inside the adult eye. It arises from pigment-producing cells in the choroid and may threaten vision or spread to distant organs.

\section*{Assessment and Management}
Assessment is guided by the history, examination, and appropriate investigations. Management should be individualised by a qualified clinician and follow current Indian clinical guidance. Depending on the condition, care may include observation, medicines, procedures, surgery, rehabilitation, or specialist referral. Self-treatment should not delay evaluation of serious symptoms.

\section*{When to Seek Medical Care}
New flashes, floaters, visual-field loss, blurred vision, or a visible change in the eye requires prompt ophthalmic assessment.

\section*{India-specific Care}
In India, start with a qualified general physician or specialist at an appropriate clinic or hospital. Emergency symptoms should be taken to the nearest emergency department. Treatment availability varies by facility, so referral to a medical college or tertiary centre may be needed for complex disease.

\section*{Sources and Further Reading}
\begin{itemize}
    \item Indian Council of Medical Research (ICMR). Clinical and public-health guidance.
    \item Ministry of Health and Family Welfare, Government of India.
    \item World Health Organization. Relevant clinical and public-health information.
\end{itemize}

